% =============================
% ABSTRACT
% =============================

\chapter*{Abstract}
This thesis examines how seriously quantum computers threaten the security of today’s encryption and to what extent data would be endangered if a market-ready quantum computer were available tomorrow. 
This thesis introduces the physical principles of quantum computing, including qubits, superposition, entanglement, and quantum gates, and contrasts their computational potential with that of classical computers. 
The current state of development and realistic timelines for fault-tolerant devices with thousands of logical qubits are outlined. 
Building on this foundation, modern cryptographic methods are explained, focusing on the role of symmetric and asymmetric schemes and the dependence of RSA on the hardness of integer factorization. 
Using brute-force attacks and Shor's algorithm as central examples, it is shown how a sufficiently powerful quantum computer could break widely used public-key systems in polynomial time.
The resulting risks for long-term confidentiality under ``harvest now, decrypt later''-attacks are assessed for governments, companies, and private individuals. An expert interview in the field of cryptography serves to validate and complement the theoretical findings.
Finally, the societal and economic impacts are discussed and the future outlook through post-quantum cryptography is evaluated.
The thesis concludes that although the day of market-ready quantum computers has not yet arrived, immediate migration strategies are essential to ensure future-proof data security.

\chapter*{Kurzfassung}
Diese Berufsmaturitätsarbeit untersucht, wie ernsthaft Quantencomputer die Sicherheit der heutigen Verschlüsselung bedrohen und inwieweit Daten gefährdet wären, wenn morgen ein marktreifer Quantencomputer verfügbar wäre.
Diese Arbeit stellt die physikalischen Prinzipien des Quantencomputings vor, darunter Qubits, Superposition, Verschränkung und Quantengatter, und vergleicht deren Rechenpotenzial mit dem klassischer Computer.
Der aktuelle Entwicklungsstand und realistische Zeitpläne für fehlertolerante Geräte mit tausenden von logischen Qubits werden skizziert. 
Auf dieser Grundlage werden moderne kryptografische Methoden erläutert, wobei der Schwerpunkt auf der Rolle symmetrischer und asymmetrischer Verfahren und der Abhängigkeit von RSA von der Schwierigkeit der Faktorisierung ganzer Zahlen liegt. 
Anhand von Brute-Force-Angriffen und Shors Algorithmus als zentrale Beispiele wird gezeigt, wie ein ausreichend leistungsfähiger Quantencomputer weit verbreitete Public-Key-Systeme in polynomieller Zeit entschlüsseln könnte. 
Die daraus resultierenden Risiken für die langfristige Vertraulichkeit bei „Harvest now, decrypt later”-Angriffen werden für Regierungen, Unternehmen und Privatpersonen bewertet. Zur Validierung und Ergänzung der theoretischen Erkenntnisse dient ein Experteninterview im Bereich Kryptografie.
Abschliessend werden die gesellschaftlichen und wirtschaftlichen Auswirkungen diskutiert und die Zukunftsaussicht durch Post-Quanten-Kryptografie bewertet. 
Die Arbeit kommt zu dem Schluss, dass der Tag der marktreifen Quantencomputer zwar noch nicht gekommen ist, aber sofortige Migrationsstrategien unerlässlich sind, um eine zukunftssichere Datensicherheit zu gewährleisten.
